\cotefile{84}
\scanpage{113}

% manuscript lines / position: 1–2
Je dis que cette formule implique que le diagramme similaire au précédent

% manuscript lines / position: 3–6
\diagram{\begin{tikzcd}[ampersand replacement=\&,column sep=large]
 A_1\otimes_k A_2 \arrow[r] \arrow[d,"{\deleted{\unsure{N_1}}\,n_1\otimes_k\mathrm{id}_{A_2}}"'] \arrow[d,phantom,"{=N_{12}}" {xshift=2em}] \& A \arrow[d,"N_2",no head]\\
 A_2 \arrow[r,shift left=.35ex,no head] \arrow[r,shift right=.35ex] \& A_2
\end{tikzcd}}

% manuscript lines / position: 4–7, left of diagram, arrow from left label
\margin{déduite de la forme $n_1$ par chgt de base $k\longrightarrow A_2$}

% manuscript lines / position: 8–10
C'est en tout cas vrai pour les éléments de \deleted{la forme} $A_1\otimes_k A_2$ de la forme $x_1\otimes x_2$, car les deux mbres sont alors

% manuscript lines / position: 11
\[
 t_1(x_1)\cdot\deleted{\unsure{\alpha(x_2)}}x_2^2
\]

% manuscript lines / position: 12–15
par quadraticité de $N_{1,2}$, $N_2$. Il reste à voir qu'on a commutativité pour les \deleted{formes} bil. associées, or pour $N_{12}$ c'est la forme déduite de

% manuscript lines / position: 16
\[
 \deleted{(x,y)\longrightarrow}\quad \varphi_{n_1}(x,y)=t_1(x\bar y)\qquad\text{sur } A_1\times A_1
\]

% manuscript lines / position: 17–18
par chgt. de base $k\longrightarrow A_1$, or on a si \inserted{$x'=\alpha_1(x),\,y'=\alpha_1(y)$}

% manuscript lines / position: 18–20
\[
 \left\{\begin{aligned}
 \deleted{\unsure{x'}}\,\alpha_1(\bar y)&=\sigma_2(\alpha_1(y))\\
 \underset{\text{dans }A_2}{t_1(x\bar y)}&=T_2(\deleted{\unsure{\alpha|\alpha_1}}(x\bar y))\; \unsure{.}
 \end{aligned}\right.
\]

% manuscript lines / position: 21–22
\[
 \text{donc}\qquad \varphi_{n_1}(x,y)=\underset{\substack{\unsure{\mathord{''}|=}\\ \varphi_{N_2}(x',y')}}{T_2(\unsure{\cdot}x'\sigma_2(y'))}\qquad \text{O.K.}
\]

% manuscript lines / position: bottom left
% archive sheet number: 112
